\documentclass[11pt,a4paper]{article}

\usepackage[T1]{fontenc}
\usepackage[utf8]{inputenc}
% newtx is installed but broken on this system (inputs a missing binhex.tex); libertinus works.
\IfFileExists{libertinus.sty}{\usepackage{libertinus}}{\usepackage{lmodern}}
\usepackage[margin=2.3cm,bottom=2.8cm]{geometry}
\usepackage{amsmath,amssymb,amsthm}
\usepackage{booktabs}
\usepackage{microtype}
\usepackage{tcolorbox}
\usepackage{enumitem}
\usepackage{graphicx}
\usepackage[hidelinks]{hyperref}
\tcbuselibrary{skins,breakable}

% ---- palette: ochre = discrete counting channel, teal = continuous market channel
\definecolor{ochre}{HTML}{A2700F}
\definecolor{teal}{HTML}{0F6D6A}
\definecolor{warn}{HTML}{9C3355}
\definecolor{ink}{HTML}{151C21}
\definecolor{rule}{HTML}{D3DBDE}

\newtcolorbox{plain}[1][]{
  breakable, enhanced, colback=teal!4, colframe=teal!55, boxrule=0.4pt,
  left=8pt,right=8pt,top=6pt,bottom=6pt, arc=1pt,
  fonttitle=\bfseries\sffamily\footnotesize, coltitle=teal,
  title={IN PLAIN ENGLISH}, #1}

\newtcolorbox{worked}[1][]{
  breakable, enhanced, colback=ochre!5, colframe=ochre!60, boxrule=0.4pt,
  left=8pt,right=8pt,top=6pt,bottom=6pt, arc=1pt,
  fonttitle=\bfseries\sffamily\footnotesize, coltitle=ochre,
  title={#1}}

\newtcolorbox{point}[1][]{
  breakable, enhanced, colback=black!3, colframe=ink!45, boxrule=0.4pt,
  left=8pt,right=8pt,top=6pt,bottom=6pt, arc=1pt,
  fonttitle=\bfseries\sffamily\footnotesize, coltitle=ink,
  title={THE POINT}, #1}

\newtcolorbox{warnbox}[1][]{
  breakable, enhanced, colback=warn!4, colframe=warn!55, boxrule=0.4pt,
  left=8pt,right=8pt,top=6pt,bottom=6pt, arc=1pt,
  fonttitle=\bfseries\sffamily\footnotesize, coltitle=warn,
  title={#1}}

\theoremstyle{plain}
\newtheorem{claim}{Claim}
\newtheorem{proposition}{Proposition}

\setlist[itemize]{leftmargin=1.3em, itemsep=2pt, topsep=3pt}
\setlist[enumerate]{leftmargin=1.6em, itemsep=3pt, topsep=3pt}

\title{\vspace{-1.2cm}\textbf{Filtering a Latent Intensity from Three Channels}\\[4pt]
\large A primer, with worked examples, on what is actually wrong\\ with piecewise Gamma--Poisson in-play updating}
\author{}
\date{}

\begin{document}
\maketitle
\vspace{-1.4cm}

\begin{center}\rule{\textwidth}{0.4pt}\end{center}
\vspace{-0.2cm}

\section*{How to read this}

The long cross-domain report is reference material. This is the entry point. It makes
\textbf{three claims}, gives each one a plain-English gloss, a formal statement, and a
worked example with real arithmetic --- and then tells you the single experiment that
decides whether any of it is worth building.

Every number in the worked examples was computed, not asserted. Where a result cuts
against the thesis, it is reported as such (see Example~3).

\begin{point}[title={THE WHOLE DOCUMENT IN SIX LINES}]
\begin{enumerate}[label=\textbf{\arabic*.}]
\item Gamma--Poisson conjugate updating is the \emph{exact} filter --- for a $\lambda$ that
      \emph{does not move}. That is the bug. It never forgets.
\item Your event timestamps identify $\lambda$ under $\mathbb{P}$. Your odds identify
      $\lambda$ under $\mathbb{Q}$. These are different numbers.
\item The odds have \emph{already seen} your events. Fusing them double-counts.
\end{enumerate}
Claim~3 raises a prior question that makes most of the rest moot, and it is cheap to answer.
Go to \S\ref{sec:experiment}.
\end{point}

\section{The setup, stripped of context}

We track a system on a finite horizon $t\in[0,T]$. A latent, time-varying intensity
$\lambda(t)>0$ governs the arrival of discrete events. We have exactly three sources of
information, and no others:

\begin{center}
\begin{tabular}{@{}llp{7.4cm}@{}}
\toprule
\textbf{Channel} & \textbf{Symbol} & \textbf{What it is} \\
\midrule
C0 & $\lambda_0 \sim \Gamma(a_0,b_0)$ & A static prior from history. \\
C1 & $N_t$, jumps at $\tau_1,\tau_2,\dots$ & Exact event timestamps. A counting process whose intensity \emph{is} the thing we want. \\
C2 & $Z_t$ & A live market price on the same outcome, from which an implied $\lambda$ can be inverted. \\
\bottomrule
\end{tabular}
\end{center}

\noindent
The object we want is the conditional law $\pi_t(\cdot) = \mathrm{Law}\big(\lambda_t \mid \mathcal{Y}_t\big)$,
where $\mathcal{Y}_t = \sigma(N_s, Z_s : s\le t)$ --- and, for pricing, the law of the
\textbf{remaining integrated intensity}
\[
  \Lambda_{t,T} \;=\; \int_t^T \lambda_s\,ds,
  \qquad\text{because}\qquad
  N_T - N_t \,\big|\, \Lambda_{t,T} \;\sim\; \mathrm{Poisson}\big(\Lambda_{t,T}\big).
\]

\begin{plain}
$\Lambda_{t,T}$ is ``how many more events do we expect, integrated over the time that is
left.'' It is the \emph{only} thing that matters for pricing. Not $\lambda$ at this instant ---
the integral of $\lambda$ over the remaining horizon. As $t\to T$ this goes to zero no matter
how large $\lambda$ is, which is why everything gets strange near expiry.
\end{plain}

\noindent
This is a textbook problem with a name: \textbf{nonlinear filtering of a doubly stochastic
(Cox) intensity under mixed observation channels}. It is not a football problem. The mature
literature on the exact three-channel version is \emph{reduced-form credit risk}, because a
credit default swap is precisely a live market price on a latent arrival intensity.

\section{Claim 1: your filter is exactly right, for the wrong model}

\begin{claim}
Gamma--Poisson conjugate updating is the exact, finite-dimensional, closed-form optimal
nonlinear filter --- for a \emph{static} Cox intensity. It is optimal and it is degenerate.
\end{claim}

Let $\lambda_t = \lambda_0\,\Phi(t)$ with $\Phi$ a \emph{known deterministic} modulation
(injury-time multipliers, score-state multipliers --- anything measurable with respect to the
observed history). By Jacod's formula the likelihood of an observed path is
\[
  \frac{d\mathbb{P}^\lambda}{d\mathbb{P}^1}\bigg|_{\mathcal{F}_t}
  \;=\; \exp\!\left(\int_0^t \log\lambda_s\,dN_s - \int_0^t(\lambda_s-1)\,ds\right)
  \;\propto\; e^{-\lambda_0\int_0^t\Phi(s)ds}\;\lambda_0^{N_t}\prod_{i}\Phi(\tau_i),
\]
so that, multiplying by the $\Gamma(a_0,b_0)$ prior,
\begin{equation}
\boxed{\;
  \lambda_0 \,\big|\, \mathcal{Y}^N_t \;\sim\;
  \Gamma\Big(\underbrace{a_0 + N_t}_{\text{shape}},\;
             \underbrace{b_0 + \textstyle\int_0^t\Phi(s)\,ds}_{\text{rate}}\Big).}
\label{eq:conj}
\end{equation}

Two things to notice. The known modulation $\Phi$ enters \emph{only} through the compensator
$\int_0^t\Phi\,ds$ --- the accumulated \emph{exposure} --- and never through the count. And
the rate parameter only ever \emph{grows}.

\begin{plain}
The rate parameter $b$ is an ``amount of evidence'' counter. It starts at $b_0$ (how much
history the prior is worth) and only ever increases. So the filter's willingness to be moved
by a new event only ever \emph{decreases}. It has no mechanism for forgetting, because it was
derived under the assumption that there is nothing to forget --- that $\lambda$ is a fixed
constant we are gradually learning. If $\lambda$ genuinely drifts during the match, this filter
will lag it, and the lag gets worse as the match goes on.
\end{plain}

\subsection*{Why this is not an abstract concern}

The sensitivity of the posterior mean to one additional event is
\[
  \frac{\partial}{\partial N_t}\,\mathbb{E}[\lambda_0\mid\mathcal{Y}_t]
  \;=\; \frac{1}{b_0 + \int_0^t \Phi\,ds}.
\]
So the \emph{entire} responsiveness of the filter is set by $b_0$: the strength of the prior.
And here is the problem.

\begin{worked}[WORKED EXAMPLE 1 --- an honest prior makes in-play updating do nothing]
Work in match-units of exposure, so a full match contributes $\int_0^1\Phi\,ds = 1$. Take a
prior mean of $1.5$ goals. Observe a full match in which the team scores $2$.

\smallskip
\begin{center}
\small
\begin{tabular}{@{}lcccc@{}}
\toprule
\textbf{Prior} & $\Gamma(a_0,b_0)$ & \textbf{Prior} $\lambda$ & \textbf{Posterior} $\lambda$ & \textbf{Move/goal} \\
\midrule
Honest (38 matches of history) & $\Gamma(57,38)$ & $1.500\pm0.199$ & $1.513\pm0.197$ & $\mathbf{0.026}$ \\
Weak (2 matches)               & $\Gamma(3,2)$   & $1.500\pm0.866$ & $1.667\pm0.745$ & $0.333$ \\
Zou et al.\ ($r_1 = E_H(45)$)  & $\Gamma(0.75,0.5)$ & $1.500\pm1.732$ & $1.833\pm1.106$ & $\mathbf{0.667}$ \\
\bottomrule
\end{tabular}
\end{center}

\smallskip
With a \textbf{statistically honest} prior --- one whose variance actually reflects a season of
data --- watching an entire match moves your estimate of the team's scoring rate by
\textbf{0.013 goals}, and shrinks your uncertainty by \textbf{1\%}. The in-play information is
completely swamped. The filter is not wrong; it is telling you the truth: \emph{two goals is
almost no evidence against 38 matches of history.}

\smallskip
This is exactly why Zou et al.\ report that the variance-matched prior gives $r_1\in(30,150)$
and ``leads to no improvement,'' and why they abandon it for $r_1 = E_H(45)\approx0.75$ --- a
prior worth \emph{half a match}. That is not a modelling choice. It is a fudge factor, chosen
to make the filter move.
\end{worked}

\begin{point}
The conjugate filter's inability to respond is not a bug in the arithmetic --- it is the
\emph{correct} answer to the wrong question. Asking ``what is this team's fixed strength?''
correctly returns ``the same as before; you have shown me almost no data.'' The question you
actually want to ask is \textbf{``has this team's strength \emph{changed}?''} --- and that
requires a model in which it \emph{can} change.
\end{point}

\subsection*{The fix: discount the evidence counter}

Let the state evolve by a multiplicative Beta shock (Smith's power-steady model;
Harvey--Fernandes local-level count model). Gamma conjugacy is preserved \emph{exactly}, and
the prediction step is a one-line discounting of the sufficient statistics:
\begin{equation}
\small
\boxed{\;
  \Gamma(a_{t-1},\,b_{t-1})
  \;\xrightarrow{\ \text{predict}\ }\;
  \Gamma(\delta a_{t-1},\,\delta b_{t-1})
  \;\xrightarrow{\ \text{update}\ }\;
  \Gamma\big(\delta a_{t-1} + \Delta N_t,\ \ \delta b_{t-1} + \Delta\!\!\int\!\Phi\big),
  \qquad \delta\in(0,1].}
\label{eq:discount}
\end{equation}
The posterior \emph{mean} is preserved by the prediction step ($\delta a/\delta b = a/b$); the
\emph{variance} is inflated by $1/\delta$. That is precisely ``I still believe the same thing,
but I am less sure of it than I was a minute ago.''

\begin{plain}
$\delta$ is a forgetting rate. Every minute, the filter throws away a fraction $(1-\delta)$ of
its accumulated evidence. This stops the evidence counter $b$ from growing without bound, so
the filter stays responsive for the whole match instead of freezing up. Setting $\delta=1$
recovers exactly what you have now.
\end{plain}

\begin{worked}[WORKED EXAMPLE 2 --- the discount factor \emph{derives} the number they guessed]
The evidence counter now obeys $b_t = \delta b_{t-1} + \Delta$, with $\Delta = 1/90$ per minute.
This is a geometric series with a \emph{fixed point}:
\[
  b^\star \;=\; \frac{\Delta}{1-\delta}
  \qquad\text{(the steady-state ``effective memory'' of the filter).}
\]

\smallskip
\begin{center}
\begin{tabular}{@{}lcc@{}}
\toprule
$\delta$ (per minute) & Effective memory $b^\star$ & Half-life \\
\midrule
$0.9800$ & $0.556$ matches & $34.3$ min \\
$\mathbf{0.98519}$ & $\mathbf{0.750}$ \textbf{matches} & $\mathbf{46.4}$ \textbf{min} \\
$0.9900$ & $1.111$ matches & $69.0$ min \\
$1.0000$ & $\infty$ (never forgets) & $\infty$ \\
\bottomrule
\end{tabular}
\end{center}

\smallskip
Solve $b^\star = 0.75$ for $\delta$ and you get $\delta = 0.98519$ --- whose forgetting
half-life is \textbf{46.4 minutes}, i.e.\ almost exactly one half of football.

\smallskip
So Zou et al.'s hand-picked $r_1 = E_H(45)$ is \emph{exactly} the steady state of a discount
filter with a one-half half-life. Their fudge factor is a real parameter in disguise. The
difference is that $\delta$ is (i) principled, (ii) applied \emph{recursively} rather than
once at kick-off, and (iii) \textbf{estimable from data} --- you can fit it by maximising the
one-step-ahead predictive likelihood over a historical panel, which you cannot do with a
hyperparameter you fixed by hand.
\end{worked}

\begin{point}
This is a \textbf{one-line change} to your existing updater, it preserves conjugacy exactly,
it costs two multiplications, and it replaces a guessed constant with a fitted one. If you do
nothing else from this document, do this.
\end{point}

\section{Claim 2: events and odds measure different things}

\begin{claim}
The counting channel identifies $\lambda$ under the physical measure $\mathbb{P}$. The market
channel identifies $\lambda$ under the pricing measure $\mathbb{Q}$. These differ by a
multiplicative premium $\gamma$, and $\gamma\ne1$.
\end{claim}

The Girsanov change of measure for a point process multiplies the intensity by a predictable,
positive process:
\begin{equation}
\boxed{\;\lambda^{\mathbb{Q}}_t \;=\; \gamma_t\cdot\lambda^{\mathbb{P}}_t\;}
\qquad
\frac{d\mathbb{Q}}{d\mathbb{P}}\bigg|_{\mathcal{F}_t}
= \exp\!\left(\int_0^t\!\log\gamma_s\,dN_s - \int_0^t\!(\gamma_s-1)\lambda^{\mathbb{P}}_s ds\right).
\label{eq:girsanov}
\end{equation}

In credit risk $\gamma$ is the \textbf{jump-to-default risk premium}, and it is measured
directly: estimate $\lambda^{\mathbb{P}}$ from realised default \emph{events}, estimate
$\lambda^{\mathbb{Q}}$ from CDS \emph{prices}, take the ratio. It comes out around $2\times$,
it is time-varying, and it is strongly counter-cyclical. In your setting $\gamma$ is the
overround plus whatever systematic distortion the market carries (a favourite--longshot bias
is exactly a $\gamma$ that varies with the level of $\lambda$).

\begin{plain}
The market price is not a noisy measurement of the true rate. It is a noisy measurement of the
\emph{price} of the true rate --- which includes the bookmaker's margin and whatever
risk-aversion or bias the market carries. If you average your event-based estimate with your
odds-based estimate as though they were two thermometers reading the same room, you get a
number that is systematically off in the same direction for the whole match. That is worse
than noise: noise averages out, bias does not, and Kelly staking on a biased estimate loses
money in a very specific and repeatable way.
\end{plain}

\begin{worked}[WORKED EXAMPLE 3 --- inverting a totals ladder \emph{(and a negative result)}]
Take a totals market: \textbf{Over 2.5 @ 1.90}, \textbf{Under 2.5 @ 2.00}. Raw implied
probabilities $1/1.90 = 0.5263$ and $1/2.00 = 0.5000$ sum to $1.0263$: an overround of
$2.63\%$.

\medskip
\textbf{Step 1 --- de-vig.} Two methods:
\begin{center}
\begin{tabular}{@{}lcc@{}}
\toprule
Method & $\mathbb{Q}(N\ge3)$ & Implied $\Lambda$ \\
\midrule
Proportional (divide by the book sum) & $0.5128$ & $2.7264$ \\
Shin (1993), insider-trading model, $z=0.0263$ & $0.5132$ & $2.7278$ \\
\midrule
\multicolumn{2}{@{}l}{\emph{Difference}} & $\mathbf{+0.0014}$ \textbf{goals} \\
\bottomrule
\end{tabular}
\end{center}

\textbf{This is a negative result and I am reporting it as one.} On a \emph{tight two-way}
book, the choice of de-vig algorithm is worth \emph{one thousandth of a goal}. It does not
matter. Do not spend a week on it.

\smallskip
It \emph{does} matter on longshots. On a $1.25\,/\,6.00\,/\,13.00$ three-way book (overround
$4.36\%$), Shin moves the favourite by $+1.41$ pp and the $13.0$ longshot by $-0.82$ pp --- an
$11\%$ \emph{relative} error on the longshot. So: use Shin or a power method on 1X2 and
correct-score, and don't agonise over totals.

\medskip
\textbf{Step 2 --- the inversion, and its sensitivity.} Invert $\mathbb{Q}(N\ge3) = 0.5128$
under $N\sim\mathrm{Poisson}(\Lambda)$ to get $\Lambda = 2.7264$. Now note the identity
\[
  \frac{\partial}{\partial\Lambda}\,\mathbb{P}(N\ge k) \;=\; \mathrm{pmf}(k-1;\Lambda),
  \qquad\text{here}\qquad
  \frac{\partial\mathbb{Q}(N\ge3)}{\partial\Lambda} = \mathrm{pmf}(2;2.726) = 0.2433.
\]
This is the ``vega'' of the contract --- how much the \emph{price} moves per unit of
\emph{parameter}. Invert it:
\[
  \boxed{\;\text{a \emph{one-tick} } (0.01) \text{ price error} \;\Longrightarrow\;
  \Delta\Lambda = 0.01/0.2433 = \mathbf{0.041}\ \textbf{goals}.\;}
\]
\emph{That} is the number that should worry you --- not the de-vig method. A single tick of
bid--ask noise moves your implied intensity by four hundredths of a goal. This is the
ill-posedness of the inverse problem, made concrete.
\end{worked}

\begin{point}
The vega gives you the observation variance \emph{for free}, and it makes the market channel
\textbf{self-weighting}:
\[
  R_t \;\asymp\;
  \left(\frac{\tfrac12\big(\text{ask}-\text{bid}\big)}{\big|\partial C/\partial\lambda\big|}\right)^{\!2}
  + R^{\text{model}}.
\]
As liquidity dries up the spread widens and $R_t$ grows. As the contract becomes insensitive to
$\lambda$ (deep in/out of the money, or $t\to T$ where $\partial C/\partial\lambda\to0$), the
denominator collapses and $R_t$ \emph{explodes}. Either way the Kalman gain on the market
channel automatically goes to zero. \textbf{You never hand-tune a weight for the market signal.
The spread-over-sensitivity ratio is the weight.}
\end{point}

\subsection*{One market cannot identify two intensities}

\begin{proposition}[Rank condition]
If the latent state is $\boldsymbol\lambda = (\lambda^{(1)},\dots,\lambda^{(d)})$ and you observe
$m$ market functionals $g_1,\dots,g_m$, local identification requires
$\operatorname{rank}\big(\partial g_i/\partial\lambda^{(j)}\big) = d$. An \emph{aggregate}
market is a function of $\sum_j\lambda^{(j)}$ alone; its Jacobian has rank $1$.
\end{proposition}

\begin{worked}[WORKED EXAMPLE 4 --- the totals market is blind to the split]
Two completely different systems:
\begin{center}
\begin{tabular}{@{}lccc@{}}
\toprule
$(\lambda^{(1)},\lambda^{(2)})$ & Sum & $\mathbb{P}(N\ge3)$ & Outcome probabilities (1/X/2) \\
\midrule
$(2.0,\;0.8)$ & $2.8$ & $\mathbf{0.5305}$ & $0.654$ / $0.205$ / $0.141$ \\
$(1.4,\;1.4)$ & $2.8$ & $\mathbf{0.5305}$ & $0.374$ / $0.253$ / $0.374$ \\
\bottomrule
\end{tabular}
\end{center}
The totals price is \textbf{identical to four decimal places} --- it must be, since the sum of
independent Poissons is Poisson in the sum. But the home-win probability differs by
\textbf{28.1 percentage points}.

\smallskip
So: a totals market supplies \emph{exactly one} linear constraint, on the sum. It contains
\emph{zero} information about the split. Inverting a totals ladder and claiming to have
recovered two team intensities is not an approximation --- it is unidentified. You need a
second, differently-weighted functional (a spread/handicap market, sensitive to
$\lambda^{(1)}-\lambda^{(2)}$) to get rank $2$.
\end{worked}

\begin{point}
Similarly, a \emph{single maturity} identifies only the \emph{level} $\Lambda_{t,T}=\int_t^T\lambda_s ds$,
never the shape of $\lambda_s$ on $(t,T)$. Any ``instantaneous implied intensity'' extracted from a
single full-match market is an artefact of your own $\Phi$. Recovering the shape needs a calendar
spread --- a market at an interim maturity --- exactly as an instantaneous forward hazard rate
requires a CDS \emph{term structure}, not a single quote.
\end{point}

\section{Claim 3: the odds have already seen your events}

\begin{claim}
Optimal filtering requires the observation channels to be conditionally independent given the
state. They are not. The market price at time $t$ is a functional of an information set that
\emph{contains} $\{N_s : s\le t\}$. A naive product-of-likelihoods fusion double-counts the
point process.
\end{claim}

This is the formal content of Duffie--Lando (2001): when the market observes the state only
through noisy information, the intensity \emph{perceived by the market} is not the structural
hazard --- it is the conditional hazard given the market's filtration. \textbf{The market's
implied intensity is itself a posterior.} It is not a sensor reading of $\lambda$; it is a
noisy reading of \emph{another agent's filter} of $\lambda$, run on a strictly larger
information set than yours.

\begin{plain}
You and the bookmaker are both watching the same match. When a goal goes in, you update, and
the odds move --- because \emph{they} updated too, on the same goal. If you then treat the new
odds as independent corroborating evidence, you have counted that goal twice. Your filter
becomes more confident than the evidence warrants, and because Kelly stakes scale
\emph{inversely with variance}, an over-confident filter systematically overstakes.
\end{plain}

\begin{worked}[WORKED EXAMPLE 5 --- what the double-count actually costs you]
Suppose the event-only filter gives $\hat x^N$ with variance $P_N = 0.09$ (sd $0.30$), the
market gives $\hat x^Z$ with variance $P_Z = 0.04$ (sd $0.20$), and their true correlation is
$\rho = 0.9$ (high, because they are largely watching the same goals).

\smallskip
The naive Kalman fusion assumes $\rho=0$, weights by inverse variance
($w_N = 0.308$, $w_Z = 0.692$), and \emph{reports}
\[
  P_{\text{naive}} = \big(P_N^{-1} + P_Z^{-1}\big)^{-1} = 0.0277 \quad (\text{sd } 0.166).
\]
But the \emph{actual} variance of that same weighted combination, with $\rho=0.9$, is
\[
  P_{\text{true}} = w_N^2 P_N + w_Z^2 P_Z + 2w_N w_Z \rho\sqrt{P_N P_Z} = 0.0507
  \quad (\text{sd } 0.225).
\]

\smallskip
\begin{center}
\begin{tabular}{@{}lcc@{}}
\toprule
& Variance & sd \\
\midrule
What the filter \emph{reports} & $0.0277$ & $0.166$ \\
What is \emph{true} & $0.0507$ & $0.225$ \\
\midrule
\textbf{Variance understated by} & \multicolumn{2}{c}{$\mathbf{83\%}$} \\
\textbf{Kelly overstakes by} & \multicolumn{2}{c}{$\mathbf{1.83\times}$} \\
\bottomrule
\end{tabular}
\end{center}

\smallskip
You would bet \textbf{nearly twice} the correct fraction of bankroll, on every bet, forever ---
and you would never see it in your backtest, because the backtest would use the same broken
variance.
\end{worked}

\subsection*{Three fixes, in increasing order of rigour}

\begin{enumerate}
\item \textbf{Innovation orthogonalisation.} Do not feed the market's \emph{level} into the
filter. Feed its \emph{residual} against your own event-only filter:
$\tilde Z_t = Z_t - \mathbb{E}[Z_t \mid \mathcal{Y}^N_t]$. Everything the market says that your
events already implied is projected out; only its genuinely private information updates the
state. Cost: run two filters.

\item \textbf{Covariance intersection} (Julier--Uhlmann 1997). The robust fusion rule when the
cross-correlation is \emph{unknown}:
\[
  \big(P^{\mathrm{CI}}\big)^{-1} = \omega P_N^{-1} + (1-\omega)P_Z^{-1},
  \qquad \omega = \arg\min \det P^{\mathrm{CI}}.
\]
It is guaranteed \emph{non-divergent for any} correlation, at the price of conservatism. In
Example~5 it returns $\omega^\star=0$, $P^{\mathrm{CI}} = 0.040$, sd $0.200$ --- which never
understates the truth. Given that our correlation is not merely unknown but \emph{known to be
high}, that conservatism is exactly what we want.

\item \textbf{Model the market as an agent.} Posit that the market runs its own filter on
$\mathcal{G}_t = \mathcal{F}^N_t \vee \sigma(U_{[0,t]})$ with $U$ its private signal, so that
$Z_t = c_t\,\hat x^N_t + (1-c_t)\,\hat U_t + \eta_t$, and \emph{solve for} $\hat U_t$ before
fusing. This is the most principled route, and it is what \S\ref{sec:experiment} estimates.
\end{enumerate}

\section{The unified architecture, in one page}

Separate what is \emph{known} from what must be \emph{filtered}:
\begin{equation}
\small
\boxed{\;
  \lambda^{(j)}_t \;=\;
  \underbrace{\hat\lambda^{(j)}_0}_{\text{static prior (C0)}}
  \;\cdot\;
  \underbrace{\Phi_j\big(t,\,S_{t^-}\big)}_{\substack{\text{KNOWN modulation:}\\
  \text{drift}\,\times\,\text{epochs}\,\times\,\text{state}}}
  \;\cdot\;
  \underbrace{\exp\big(X^{(j)}_t\big)}_{\substack{\text{LATENT deviation}\\ \text{(the only filtered part)}}}}
\label{eq:decomp}
\end{equation}

The injury-time multipliers $\rho(t)$ and the score-state multipliers $\tau_{xy}$ are
\emph{predictable given the observed history}. They belong in $\Phi$, where they cost nothing.
Putting them in the filter would cost dimensions for no information. Only $X$ --- the log-deviation
of the true intensity from its prior --- is latent.

\medskip\noindent\textbf{State dynamics} (Black--Karasinski / exponential-OU, so $\lambda>0$ automatically):
\begin{equation}
  dX^{(j)}_t \;=\;
  \underbrace{-\kappa_j X^{(j)}_t\,dt}_{\text{mean-revert to the prior}}
  \;+\; \underbrace{\sigma_j(t)\,dW^{(j)}_t}_{\text{true strength drifts}}
  \;+\; \underbrace{\textstyle\sum_k \zeta_{j,k}\,dH^{(k)}_t}_{\text{structural breaks (red cards)}},
  \qquad
  d\beta_t = -\kappa_\beta(\beta_t - \bar\beta)dt + \sigma_\beta dW^\beta_t,
\label{eq:sde}
\end{equation}
where $\beta = \log\gamma$ is the \textbf{market bias state} from Claim~2. Note
$\kappa$ is the continuous-time forgetting rate: $\delta \approx e^{-\kappa\Delta}$. It is the
same parameter as \eqref{eq:discount}, and it is the single most important number in the system.

\medskip\noindent\textbf{The filter} (Kushner--Stratonovich with mixed channels):
\begin{equation}
\small
\boxed{
\begin{aligned}
d\pi_t(\varphi) \;=\;& \pi_t(\mathcal{A}_t\varphi)\,dt
&&\text{\scriptsize (i) drift toward prior}\\
&+\; \textstyle\sum_m \Big[\pi_t(\varphi g_m) - \pi_t(\varphi)\pi_t(g_m)\Big]
   \big(R^{(m)}_t\big)^{-1}\,\omega^{(m)}_t\,
   \Big(dZ^{(m)}_t - \pi_t(g_m)\,dt\Big)
&&\text{\scriptsize (ii) market innovation}\\
&+\; \textstyle\sum_j \left[\frac{\pi_{t^-}(\varphi\lambda^{(j)})}{\pi_{t^-}(\lambda^{(j)})} - \pi_{t^-}(\varphi)\right]
   \Big(dN^{(j)}_t - \pi_{t^-}(\lambda^{(j)})\,dt\Big)
&&\text{\scriptsize (iii) point-process innovation}
\end{aligned}}
\label{eq:ks}
\end{equation}

\begin{plain}
Read the three lines as three forces acting on your belief:
\begin{itemize}
\item \textbf{(i)} With no new data, your belief relaxes back toward the pre-match prior and
grows more uncertain. The prior is not a one-shot starting value --- it is an \emph{attractor}.
\item \textbf{(ii)} When the price moves, correct in proportion to how surprising the move was,
scaled by how noisy the quote is ($R_t$, from the spread-over-vega) and tempered by $\omega_t$
to avoid double-counting (Claim~3).
\item \textbf{(iii)} When an event fires, multiply your posterior by $\lambda$ and renormalise
(for a Gamma: $a \mapsto a+1$). When an event \emph{fails} to fire, the compensator term
$-\pi(\lambda)\,dt$ drags your belief \emph{down}, at exactly the rate you expected events to
arrive. \textbf{No news is news.}
\end{itemize}
\end{plain}

\subsection*{EKF or particle filter?}

\begin{worked}[WORKED EXAMPLE 6 --- why the Gaussian approximation fails \emph{here specifically}]
Suppose it is half-time and no goals have been scored. With prior $\Gamma(3,2)$ the exact
posterior is $\Gamma(3,\,2.5)$: mean $1.200$, sd $0.693$, \textbf{skewness $1.155$}.

The EKF replaces this with $\mathcal{N}(1.200,\,0.693^2)$. That Gaussian:
\begin{itemize}
\item assigns $\mathbb{P}(\lambda<0) = \mathbf{4.2\%}$ --- to a rate parameter;
\item and, far more damagingly, \emph{mis-states the upper tail}:
\end{itemize}
\begin{center}
\begin{tabular}{@{}lccc@{}}
\toprule
& True ($\Gamma$) & Gaussian (EKF) & Error \\
\midrule
$\mathbb{P}(\lambda > 2.5)$ & $0.0517$ & $0.0303$ & $\mathbf{-41\%}$ \\
$\mathbb{P}(\lambda > 3.0)$ & $0.0203$ & $0.0047$ & $\mathbf{-77\%}$ \\
\bottomrule
\end{tabular}
\end{center}

\smallskip
And the tail is \emph{the thing you are pricing}: an over/under contract is
$\mathbb{P}(N_T - N_t \ge k)$, a tail functional. The EKF is wrong by a factor of four exactly
where the money is.
\end{worked}

\begin{center}
\small
\begin{tabular}{@{}p{3.5cm}p{5.1cm}p{5.9cm}@{}}
\toprule
& \textbf{Continuous-discrete EKF} & \textbf{Sequential Monte Carlo} \\
\midrule
Cost per step & $O(d^2)$, $d\approx2$ --- negligible & $O(Pd)$, $P\sim10^3$--$10^4$ --- \emph{also} negligible \\
Survival term $e^{-\int\lambda}$ & Gaussian moment-matching only & Exact (a per-particle path integral) \\
Low-count skew & \textbf{Badly wrong} (Example 6) & Exact \\
Tail functionals & Poor --- and this is what you price & Exact \\
Regime jumps in $\Phi$ & Needs re-linearisation & Trivial \\
Degeneracy risk & None & Yes --- resample when $\mathrm{ESS}<P/2$ \\
\bottomrule
\end{tabular}
\end{center}

\begin{point}
\textbf{Use SMC.} The state dimension is 1--3, the horizon is short, the events number
$\le10$: a $10^4$-particle bootstrap filter costs microseconds. The \emph{entire} argument for
the EKF is computational, and it is worth nothing here --- while the argument against it (bias
in the low-count, skewed, tail-functional regime) bites maximally. Rao-Blackwellise: conditional
on the log-OU path, the level retains Gamma conjugacy, so sample only the dynamic component and
integrate the level analytically.

\smallskip
Keep the closed-form discount filter \eqref{eq:discount} as the \emph{production} path
(nanoseconds, no tuning beyond $\delta$), and SMC as the reference against which you measure its
approximation error.
\end{point}

\section{What to actually do}
\label{sec:experiment}

\begin{warnbox}[title={THE ONE EXPERIMENT THAT DECIDES EVERYTHING}]
Claim 3 raises a prior question: \textbf{does the market channel contain anything your event
stream does not?} If the market's implied intensity is just a mechanical restatement of the
scoreline, then there is no second channel --- there is one channel and a de-vigging model, and
the entire fusion architecture is dead weight.

\smallskip
This is cheap to find out and you already have the data.
\end{warnbox}

Over your historical in-play panel, at each observation time $t$:

\begin{enumerate}
\item Run your existing conjugate filter on the \textbf{event stream alone} $\Rightarrow$
      event-only log-intensity $\hat x^N_t$.
\item De-vig the odds and invert the totals ladder $\Rightarrow$ market-implied log-intensity
      $\hat x^M_t$.
\item Regress $\hat x^M_t$ on $\hat x^N_t$ (and time-to-go). The $R^2$ measures how much of the
      market is simply \textbf{the scoreline echoed back at you}. Keep the residual $u_t$.
\item \textbf{The real test.} Regress the \emph{realised} remaining goals $N_T - N_t$ on
      $\hat x^N_t$ \textbf{and} $u_t$:
      \[
        \big(N_T - N_t\big) \;=\; \alpha + \beta_1\,\hat x^N_t + \beta_2\, u_t + \varepsilon.
      \]
\end{enumerate}

\begin{center}
\begin{tabular}{@{}lp{9.6cm}@{}}
\toprule
\textbf{Outcome} & \textbf{What it means} \\
\midrule
$\beta_2 \approx 0$ &
The market has \emph{no} incremental information beyond your events. Kill the fusion
architecture. Implement \eqref{eq:discount}, fix the de-vig, and move on. \\
\addlinespace
$\beta_2 \ne 0$ &
\textbf{That coefficient is your edge.} And it directly calibrates $\omega_t$ --- the temper in
\eqref{eq:ks} --- which you would otherwise have had to guess. \\
\bottomrule
\end{tabular}
\end{center}

\noindent
It is a day's work against the existing \texttt{DataStore}, and it determines whether the rest
of this is worth building. \textbf{Run it before you write any filter code.}

\subsection*{If it passes, the build order is}

\begin{enumerate}
\item \textbf{Discount filter} \eqref{eq:discount}. Fit $\delta$ by maximising one-step-ahead
      predictive likelihood. One line; strictly dominates what you have.
\item \textbf{Move $\Phi$ out of the filter} \eqref{eq:decomp}. Epoch and score-state multipliers
      are predictable; they are free.
\item \textbf{Self-weighting market channel} via spread-over-vega ($R_t$). No hand-tuned weights.
\item \textbf{Premium state} $\beta = \log\gamma$ \eqref{eq:girsanov}. Estimate offline first;
      if it is stable, fix it.
\item \textbf{Temper / orthogonalise} the market channel ($\omega_t$ from step 4's $\beta_2$).
\item \textbf{SMC} as reference; keep the closed-form filter in production.
\end{enumerate}

\vspace{0.4cm}
\begin{center}\rule{\textwidth}{0.4pt}\end{center}

\section*{Sources for the load-bearing claims}
\footnotesize
\sloppy
All identifiers below were retrieved from OpenAlex/arXiv and checked to resolve.

\begin{itemize}[leftmargin=1em]
\item \textbf{Snyder (1972)}, ``Filtering and detection for doubly stochastic Poisson processes,''
      \emph{IEEE Trans.\ Inf.\ Theory}. \texttt{10.1109/tit.1972.1054756} --- origin of the DSPP filter.
\item \textbf{Segall \& Kailath (1975)}, ``The modeling of randomly modulated jump processes.''
      \texttt{10.1109/tit.1975.1055359} --- the martingale-innovations form of \eqref{eq:ks}.
\item \textbf{Harvey \& Fernandes (1989)}, ``Time series models for count or qualitative observations,''
      \emph{JBES}. \texttt{10.1080/07350015.1989.10509750} --- \textbf{the discount filter \eqref{eq:discount}}.
\item \textbf{Smith (1979)}, ``A generalization of the Bayesian steady forecasting model,''
      \emph{JRSS-B}. \texttt{10.1111/j.2517-6161.1979.tb01092.x} --- the power-steady model.
\item \textbf{Lando (1998)}, ``On Cox processes and credit risky securities.'' OpenAlex \texttt{W2084716442}.
\item \textbf{Duffie \& Singleton (1999)}, ``Modeling term structures of defaultable bonds,''
      \emph{RFS}. OpenAlex \texttt{W2011886674} --- intensity-as-spread.
\item \textbf{Duffie \& Lando (2001)}, ``Term structures of credit spreads with incomplete accounting
      information,'' \emph{Econometrica}. OpenAlex \texttt{W2118891973} --- \textbf{the market's intensity is
      itself a posterior} (Claim 3).
\item \textbf{Berndt et al.}, ``Measuring default risk premia from default swap rates and EDFs.''
      \texttt{10.2139/ssrn.556080} --- direct measurement of the $\gamma$ wedge (Claim 2).
\item \textbf{Shin (1993)}, ``Prices of state-contingent claims with insider traders, and the
      favourite--longshot bias,'' \emph{Economic Journal}. \texttt{10.2307/2234526}.
\item \textbf{Breeden \& Litzenberger (1978)}, ``Prices of state-contingent claims implicit in option
      prices,'' \emph{J.\ Business}. \texttt{10.1086/296025}.
\item \textbf{Julier \& Uhlmann (1997)}, ``A non-divergent estimation algorithm in the presence of
      unknown correlations,'' \emph{ACC}. \texttt{10.1109/acc.1997.609105} --- covariance intersection.
\item \textbf{Cori et al.\ (2013)}, ``A new framework and software to estimate time-varying reproduction
      numbers during epidemics,'' \emph{Am.\ J.\ Epidemiol}. \texttt{10.1093/aje/kwt133} --- the same
      Gamma--Poisson filter, with a forgetting window, in epidemiology.
\item \textbf{Brown, Frank \& Eden (2004)}, ``Dynamic analysis of neural encoding by point process adaptive
      filtering,'' \emph{Neural Computation}. \texttt{10.1162/089976604773135069}.
\end{itemize}

\fussy
\normalsize
\vspace{0.3cm}
\noindent\emph{Caveat.} All derivations here are assembled from the cited frameworks; the specific
composition is not itself a retrieved result. The worked examples were computed numerically. A
literature search found \textbf{no} prior work fusing a historical prior, an event stream, and a
concurrent market price of the outcome outside of finance --- so the credit-risk literature is the
only mature guide to channels C1 and C2 \emph{together}.

\end{document}
